\clearpage
\VerdierPhysicalPageAnchor{0183}%
\begingroup\small\setlength{\parskip}{0pt}\setlength{\emergencystretch}{1.5em}
\setlength{\abovedisplayskip}{4pt plus 2pt minus 2pt}%
\setlength{\belowdisplayskip}{4pt plus 2pt minus 2pt}%

\VerdierSemanticAnchor{ch3d-1-3-2-morphism-m-cont}
The morphism $m$ is a monomorphism, and there is an exact sequence:

\VerdierSemanticAnchor{ch3d-1-3-2-exact-sequence}
\begin{equation}
0\longrightarrow c(\operatorname{id}_X)
\xrightarrow{m}c(\mu(S))[-1]
\longrightarrow c(\operatorname{id}_Z)[-1]
\longrightarrow0.
\tag{1.3.2.3}
\end{equation}
The complexes $c(\operatorname{id}_X)$ and $c(\operatorname{id}_Z)$ are
homotopic to zero. Consequently, the complex $c(\mu(S))$ is acyclic. The
morphism $\mu(S)$ is therefore a quasi-isomorphism.

\VerdierSemanticAnchor{ch3d-1-3-2-k-diagram}
Passing to the category $K(\mathcal A)$, we obtain a diagram (Chap. I,
3.2.2.7):
\begin{center}
\resizebox{0.84\textwidth}{!}{%
$\begin{tikzcd}[ampersand replacement=\&,column sep=3.2em,row sep=2.6em]
X \arrow[r,"\dot u"]
  \& Y \arrow[r,"\dot{p}(u)"] \arrow[dr,"\dot v"']
  \& c(\dot u) \arrow[r,"\dot{q}(u)"] \arrow[d,"\dot\mu(S)"]
  \& X[1] \\
\& \& Z \&
\end{tikzcd}$%
}
\end{center}
We then set:

\VerdierSemanticAnchor{ch3d-1-3-2-connecting}
\begin{equation}
\delta(S)=[q(u)][\mu(S)]^{-1}.
\tag{1.3.2.4}
\end{equation}
It follows from the definitions of the triangulated structures on
$K(\mathcal A)$ (Chap. I, 3.3.1) and $D(\mathcal A)$ (Chap. II, 2.2.6)
that the triangle $(X,Y,Z,[u],[v],\delta(S))$ is distinguished in
$D(\mathcal A)$. Moreover, one checks immediately that the morphism
$\delta(S)$ is functorial in $S$ ((1.3.1), property a)). We have therefore
endowed the canonical functor
$\operatorname{comp}(\mathcal A)\longrightarrow D(\mathcal A)$ with a
$\delta$-functor structure. The resulting $\delta$-functor is called the
\emph{canonical $\delta$-functor}. The morphism $\delta(S)$ is called the
\emph{connecting morphism} associated with the exact sequence $S$.

\VerdierSemanticAnchor{ch3d-1-3-3}
\medskip\noindent\textbf{1.3.3.}\quad
Let $\mathcal A^o$ be the category opposite to $\mathcal A$. For every exact
sequence $S$ in $\operatorname{comp}(\mathcal A^o)$, denote the associated
connecting morphism by $\delta^o(S)$. Let $C$ be a category, let $X$ be an
object of $C$, and let $u$ be a morphism of $C$. In the following
proposition, by abuse of notation, we denote by $X^o$ and $u^o$ the
corresponding object and morphism of the opposite category of $C$ (recall
that the opposite category of $C$ has the same sets of objects and
morphisms as $C$).

\VerdierSemanticAnchor{ch3d-1-3-4}
\medskip\noindent\textbf{Proposition 1.3.4.}\quad
\textit{a) For every exact sequence $S$ in
$\operatorname{comp}(\mathcal A)$, one has:}
\[
\delta(S[1])=-\delta(S)[1].
\]
\noindent\textit{b) For every exact sequence $S$ in
$\operatorname{comp}(\mathcal A)$, taking the identification (1.2.8) into
account, one has:}
\[
\delta^o(S^o)=\delta(S)^o[1].
\]
\endgroup
